\subsection{Model Selection}\label{sec:model-selection}

Configuration 8 is selected for the generation of price paths,  together with the $\mathrm{PnL}^{*}_{a}$ cost function branch. As listed in Tab.~\ref{tab:configurations}, configuration 8 uses a condition window of $l=10$, a noise dimension of $N=32$, generator and discriminator hidden sizes of $R_D=R_G=64$ and $n_{\mathrm{grad}} + n_{\mathrm{epochs}} = 100 + 500$. 

\textbf{Network capacity.} Configuration 9 is closer to the test split than configuration 8 on the excess kurtosis and the skewness. Its skewness seed range is also the only one that does not cross zero. However, the seed spread of configuration 9 on the excess kurtosis is more than three times as wide as the seed spread of configuration 8. Its maximum exceeds the test split, while the seed range of configuration 8 contains it (Sec.~\ref{sec:network-capacity}). Reproducibility across seeds is taken as the selection criterion, not proximity of the mean. As a result, configuration 8 is chosen over configuration 9.

\textbf{Cost function.} No branch is separable from ForGAN, so the separability criterion cannot distinguish between them. Selection among the branches is based on closeness to the test split. As reported in Sec.~\ref{sec:effect-cost}, under configuration 8 the seed ranges of the $\mathrm{PnL}^{*}_{a}$ and $\mathrm{PnL}^{*}_{a}$ \& $\mathrm{MSE}$ branches contain the test split on the skewness, the excess kurtosis, the standard deviation and the \SI{1}{\percent} quantile. The seed range of the $\mathrm{PnL}^{*}_{a}$ branch contains the maximum of the test split. The $\mathrm{PnL}^{*}_{a}$ \& $\mathrm{MSE}$ branch exceeds it across all seeds. The seed range of the $\mathrm{PnL}^{*}_{a}$ branch is closer to the test split on the standard deviation than the ForGAN baseline, and its skewness seed range is also narrower. The $\mathrm{PnL}^{*}_{a}$ branch is chosen over the ForGAN baseline and the $\mathrm{PnL}^{*}_{a}$ \& $\mathrm{MSE}$ branch. The selected model generates the price paths of Sec.~\ref{sec:option-prices}.

% ### Closing paragraph — what the selected model does not reproduce
%
% - [ ] PURPOSE: forward-looking, not a summary. Tells the reader how to
%       read Sec. 5.7. Do not restate the two decisions
%
% - [ ] Numbers are the PnL*a column of Tab. 12, NOT ForGAN. ForGAN sits
%       immediately to its left, easy to grab the wrong one
%
% - [ ] S1. The body of the distribution is close: the seed ranges contain
%       the test split on the skewness, the excess kurtosis and the
%       standard deviation. One clause, cite Sec. 5.5, no numbers
%
% - [ ] S2. The extremes do not match. Minimum -0.2542 against the
%       -0.1749 of the test split; maximum 0.3950 against 0.3229.
%       Both overshoot
%
% - [ ] S3. The autocorrelation of squared returns does not decay:
%       C2(10) is 0.153 against the 0.011 of the test split
%
% - [ ] S4. Name the consequence for Sec. 5.7 WITHOUT stating a direction:
%       the far strikes depend on the extremes, so they are affected.
%       DO NOT say whether prices come out too high or too low — that is
%       5.7's result
%
% - [ ] Keep to four sentences. This paragraph frames 5.7, it does not
%       pre-empt it
